% Introduction. Part of the team paper; included by main.tex.
\section{Introduction}
\label{sec:intro}

Large language models can perform a new task from a few examples in their prompt, with no weight update~\cite{brown2020language}.
This \emph{in-context learning} (ICL) is hard to study in natural language, because neither the task family nor the optimal
predictor is known. Garg et al.~\cite{garg2022can} made it measurable. They trained a transformer from scratch on prompts
$(x_1, f(x_1), \dots, x_n, f(x_n), x_q)$, with $f$ drawn from a known function class such as linear functions, and compared its
prediction of $f(x_q)$ with the optimal estimator. That setup gives full control over the task distribution, the number of
examples and the test-time shift.

We ask whether the same recipe works for \emph{control}. The analog of a function is a policy. The analog of a prompt is a set of
demonstrations of a task, and the query is the robot's current state. Two features make control harder than regression. First,
the task variable is not written in the demonstrations: it has to be inferred from how the demonstrations behave. Second,
predictions feed back into the inputs. A small error in one action changes the next state, so errors compound over an
episode~\cite{ross2010efficient}, and a low teacher-forced loss need not give a working controller.

We study this in a controlled setting: Push-T~\cite{florence2021implicit} in ManiSkill~\cite{tao2024maniskill3}, where the goal
position of the T varies between episodes and is \emph{hidden} from the learner. Two goal-conditioned teachers see the goal: a PPO
policy and a behavior-cloning (BC) policy distilled from it. A GPT-2-style transformer sees only demonstrations from these
teachers and the current state with the goal removed, so the demonstrations are its only source of information about the goal.
We train it on goals in three quadrants of the goal region and test it on new goals in those quadrants and in the held-out
fourth quadrant. To check that success comes from reading the context, every evaluation is repeated with demonstrations of a
\emph{different} goal and with no demonstrations, from the same initial states.

\paragraph{Contributions.}
\begin{itemize}
  \item \textbf{A controlled testbed for in-context imitation in closed-loop control.} Variable-goal Push-T with a hidden goal,
    goal-conditioned PPO and distilled BC teachers, Garg-style demonstration prompts, and correct-goal, wrong-goal and no-context
    controls on identical initial states (Section~\ref{sec:method}). Code and all result files are released, and every figure and
    table can be regenerated from them.
  \item \textbf{In-context control matches the teachers in distribution.} On unseen goals in the training quadrants, the transformer
    reaches 91--92\% success with demonstrations of the right goal, about the same as the teachers that see the goal. With wrong-goal
    or no demonstrations it reaches 1--2\%. A single demonstration already gives about 90\%.
  \item \textbf{On held-out goals, reaching transfers but holding does not.} In the held-out quadrant the transformer still reaches
    the goal in 86--89\% of episodes, but keeps the T there at the end in only 64--67\%, against 93--94\% for the teachers.
  \item \textbf{A closed-loop threshold on offline error, replicated across implementations.} A model that already prefers
    correct demonstrations offline can still fail completely in closed loop. An independent implementation reproduces every
    qualitative finding.
\end{itemize}

\paragraph{Scope.} The team's plan covers a $2 \times 2$ grid of teachers: full-state or camera observations, crossed with PPO or BC
training. This draft reports the full-state row only. The camera-based teachers are in progress and not part of these results.
%CHECK: adjust to match the proposal's framing and the camera teachers' status
